\chapter{Implementation}

% ------------------------------------------------------------
\chapter{Implementation}
% ------------------------------------------------------------
 
\section*{Introduction}
\addcontentsline{toc}{section}{Introduction}
 
Chapter~3 established the architectural design of the system —
defining what each component should do and how the layers relate to
one another. This chapter documents how that design was realised in
practice. The focus shifts from \emph{what was intended} to
\emph{what was built}: the specific hardware, software tools,
libraries, and engineering decisions that translate the architectural
model into a functioning prototype.
 
The implementation encompasses the following primary components:
 
\begin{itemize}
    \item The development environment, including hardware platform
          selection, operating system, programming tools, and
          dependency management.
    \item The embedded inference pipeline, covering model format
          conversion and the deployment challenges encountered on
          the target hardware.
    \item The Raspberry Pi software stack, including the multi-threaded
          service architecture, the Finite State Machine runtime, and
          the camera acquisition pipeline.
    \item The Arduino firmware, covering the command processor,
          actuator abstraction, and failsafe mechanism.
    \item The communication layer, including the WebSocket client
          implementation and the UART serial protocol.
\end{itemize}
 
A recurring theme throughout this chapter is the gap between
theoretical deployment assumptions and practical hardware constraints.
Several implementation decisions — most notably the choice of
inference runtime and model format — were driven not by design
preference but by the limitations of the available embedded hardware
and the compatibility boundaries of the available software toolchains.
These decisions are documented candidly so that future work can build
on the lessons learned.
 
% ------------------------------------------------------------
\section{Development Environment}
\label{sec:devenv}
% ------------------------------------------------------------
 
This section documents the hardware platform, software environment,
and dependency management process used throughout the implementation
phase. Establishing a reproducible and stable development environment
was itself a significant engineering effort, particularly with respect
to AI inference dependencies on constrained embedded hardware — a
challenge that ultimately influenced core architectural decisions and
is examined in detail in Section~\ref{sec:dep_challenges}.
 
% ............................................................
\subsection{Hardware Platform}
% ............................................................
 
The embedded hardware used in this project consists of three physical
components, each serving a distinct role in the prototype system.
 
\paragraph{Raspberry Pi 3 Model B+.}
The Raspberry Pi 3B+ is the central processing unit of the embedded
system. It features a 1.4\,GHz quad-core ARM Cortex-A53 processor,
1\,GB of LPDDR2 RAM, integrated 802.11ac dual-band Wi-Fi, and a
dedicated CSI camera port~\cite{upton2016raspberry}. While the
Cortex-A53 is a capable general-purpose processor, it lacks a
dedicated neural processing unit (NPU) or hardware floating-point
accelerator optimised for inference workloads — a limitation that
proved consequential during the model deployment phase (see
Section~\ref{sec:dep_challenges}).
 
\paragraph{Arduino Uno.}
The Arduino Uno is based on the ATmega328P 8-bit microcontroller
running at 16\,MHz with 2\,KB of SRAM and 32\,KB of flash
storage~\cite{monk2016electronics}. It is responsible exclusively for
low-level actuator control: receiving UART commands from the Raspberry
Pi and generating the PWM signals required by motor drivers and relay
circuits. Its limited resources are well-matched to this role — the
firmware is small and its execution timing is deterministic.
 
\paragraph{Raspberry Pi Camera Module 3.}
The Camera Module 3 connects to the Raspberry Pi via the CSI ribbon
interface and provides up to 12\,MP resolution through a Sony IMX708
sensor~\cite{picamera2023}. For inference purposes frames are captured
at a reduced resolution to balance image quality against the
computational cost of preprocessing and model inference on the
embedded processor.
 
% ............................................................
\subsection{Software Environment}
% ............................................................
 
\paragraph{Raspbian GNU/Linux 13 (Trixie).}
The operating system is Raspbian (Raspberry Pi OS) based on Debian 13
``Trixie'', a 64-bit Linux distribution maintained by the Raspberry Pi
Foundation and optimised for the Raspberry Pi hardware
family~\cite{raspios2024}. The 64-bit userspace is important for AI
inference libraries, many of which distribute pre-built wheels only
for \texttt{aarch64} and not for the legacy 32-bit \texttt{armv7l}
architecture.
 
\paragraph{Python 3.13.5.}
Python 3.13.5 is the primary application programming language for all
Raspberry Pi software. Python was chosen for its extensive ecosystem
of computer vision and machine learning libraries, its compatibility
with the inference runtimes under consideration, and its suitability
for rapid prototyping of complex, multi-threaded
systems~\cite{python313}.
 
\paragraph{Arduino IDE.}
The Arduino Integrated Development Environment is used to write,
compile, and flash the C++ firmware onto the Arduino Uno. The IDE
provides the AVR-GCC toolchain and the \texttt{avrdude} programmer
required to deploy firmware over USB without additional
tools~\cite{arduinoide}.
 
\paragraph{Raspberry Pi Connect.}
Raspberry Pi Connect is a secure remote access service provided by
the Raspberry Pi Foundation that enables shell and desktop access to
a Raspberry Pi over the internet without configuring port forwarding
or a VPN~\cite{rpconnect2024}. It was used during development to
remotely monitor, debug, and redeploy software on the embedded device
without requiring physical access.
 
\paragraph{OpenCV.}
OpenCV (Open Source Computer Vision Library) is an open-source
library providing optimised implementations of image processing
algorithms~\cite{bradski2000opencv}. In this project it is used for
frame preprocessing (resizing, colour space conversion, normalisation)
before frames are passed to the inference pipeline.
 
\paragraph{ONNX Runtime.}
ONNX Runtime is a cross-platform inference engine for models in the
Open Neural Network Exchange format~\cite{onnxruntime2023}. It was
selected as the primary inference backend for the server-side
deployment path due to its broad operator support, active maintenance,
and strong performance on Linux x86-64 hardware.
 
\paragraph{TensorFlow Lite.}
TensorFlow Lite is Google's inference framework for mobile and
embedded devices~\cite{david2021tensorflow}. It was evaluated as an
alternative embedded inference backend during the dependency
investigation described in Section~\ref{sec:dep_challenges}.
 
% ............................................................
\subsection{Build and Dependency Management}
% ............................................................
 
Python dependencies for the Raspberry Pi software stack are managed
using a virtual environment (\texttt{venv}) to isolate project
packages from the system Python installation. CMake is used where
native extension modules require compilation from source — a
requirement that arose during several inference library installation
attempts on the ARM platform.
 
\subsubsection{Dependency and Deployment Challenges}
\label{sec:dep_challenges}
 
One of the most time-consuming aspects of the implementation phase was
resolving inference runtime dependencies on the Raspberry Pi 3B+.
This section documents the investigation process and the decisions it
produced, because these constraints directly shaped the final
deployment architecture.
 
\paragraph{Stage 1 — Direct RF-DETR Inference.}
The initial intention was to run the RF-DETR Nano model directly on
the Raspberry Pi using its native PyTorch weights. Profiling quickly
revealed that the model was too computationally demanding for the
Cortex-A53 processor: inference times were far outside the range
acceptable for real-time operation, and peak memory usage approached
the device's 1\,GB RAM limit. A lighter-weight, hardware-optimised
format was required.
 
\paragraph{Stage 2 — ONNX Runtime on the Raspberry Pi.}
The model was exported to ONNX format (as described in
Section~\ref{sec:quantization}) and ONNX Runtime was identified as
the preferred embedded inference engine due to its Linux compatibility
and active ARM support. However, after export, no pre-built
\texttt{.whl} package for ONNX Runtime was available that matched
the combination of Python 3.13.5 and the \texttt{aarch64} architecture
of the Raspberry Pi 3B+. Attempts to install wheels built for adjacent
Python versions (3.11, 3.12) failed due to ABI incompatibilities.
Compiling ONNX Runtime from source requires CMake, a C++17-capable
compiler, and a build process that exceeds the memory capacity of the
device itself.
 
Downgrading the Python runtime to a version with available ONNX
Runtime wheels was explored, but this introduced conflicts with other
project dependencies that required Python~3.13 APIs, making this path
unviable.
 
\paragraph{Stage 3 — TensorFlow Lite via ONNX Conversion.}
TensorFlow Lite was considered as an alternative embedded runtime,
as pre-built TFLite interpreter packages are available for a wider
range of ARM Python configurations. Because RF-DETR has no native
TFLite export path, a two-stage conversion pipeline was required:
 
\begin{center}
RF-DETR (PyTorch) $\xrightarrow{\text{torch.onnx.export}}$
ONNX $\xrightarrow{\text{onnx2tf / ai\_edge\_torch}}$ TFLite
\end{center}
 
A compatible Python runtime version was identified under which the
TFLite interpreter package installed successfully. However, this
runtime version was incompatible with the system-installed OpenCV
and \texttt{picamera2} bindings, which are compiled against the
system Python and cannot be imported from a different interpreter.
Custom wrapper solutions were prototyped to bridge frame acquisition
from \texttt{picamera2} into the alternative runtime environment,
but these introduced additional fragility and maintenance overhead.
 
After completing the conversion pipeline, the resulting
\texttt{.tflite} model failed to load on the target runtime. The
RF-DETR architecture relies on custom attention and decoder operations
that are not part of the standard TFLite operator set; the conversion
tools emitted these as \emph{unsupported operators}, which the TFLite
interpreter rejected at model initialisation time. Multiple conversion
configurations and operator compatibility flags were tested without
success.
 
\paragraph{Resolution — Remote Inference Fallback.}
Having exhausted the available embedded runtime options for the current
hardware, the decision was made to adopt the remote inference fallback
path described in Section~\ref{sec:inference_pipeline}: frames are
streamed from the Raspberry Pi to the backend server over WebSocket,
where ONNX Runtime on the x86-64 server executes inference and returns
detection results to the Pi. This path has no operator-support
constraints and imposes no memory or compute requirements on the
embedded device itself.
 
This outcome reinforces the value of designing the inference pipeline
with a fallback mode from the outset (Section~\ref{sec:inference_pipeline}):
the architectural decision to separate the inference interface from its
implementation meant that switching from local to remote execution
required no changes to the surrounding software stack.
 
Table~\ref{tab:inference_stages} summarises the three stages of the
investigation and their outcomes.
 
\begin{table}[h]
\centering
\caption{Summary of embedded inference deployment investigation.}
\label{tab:inference_stages}
\begin{tabular}{p{3cm} p{4.5cm} p{4.5cm}}
\hline
\textbf{Approach} & \textbf{Blocker} & \textbf{Outcome} \\
\hline
Direct PyTorch (RF-DETR native) &
    Compute and memory exceed Pi 3B+ capacity &
    Abandoned \\
ONNX Runtime on Pi &
    No compatible wheel for Python 3.13 / \texttt{aarch64}; source build exceeds device RAM &
    Abandoned \\
TFLite via ONNX conversion &
    RF-DETR custom operators unsupported by TFLite runtime; Python runtime conflict with system OpenCV &
    Abandoned \\
Remote inference (server ONNX Runtime) &
    None &
    \textbf{Adopted} \\
\hline
\end{tabular}
\end{table}